%------------------------------ Ejercicio 5.31 -------------------------------
\setcounter{neje}{30}
\begin{nejer}
Sea $X \coloneq \mathbb{R}$, denote por $\mathcal{B} \subseteq 2^X$ a la $\sigma$-álgebra de Borel, sea $\mu \colon \mathcal{B} \to [0, \infty]$ la restricción de la medida de Lebesgue a $\mathcal{B}$, y sea $\nu \colon \mathcal{B} \to [0, \infty]$ la medida de conteo. Pruebe lo siguiente.
\begin{enumerate}[label = (\roman*), leftmargin=1.8em]
\item $\mu \ll \nu$.
\item $\mu$ no es regular interior con respecto a $\nu$.
\item No existe ninguna función medible $f \colon X \to [0, \infty]$ tal que
$$
\mu(B) = \int_B f \dd{\nu} \quad \text{para todo } B \en \mathcal{B}.
$$
\end{enumerate}
\end{nejer}
\begin{sol}
\begin{enumerate}[label = (\roman*), leftmargin=1.8em]
%----------------------------------------------------
\item % (i)
Dado $B\en\mathcal{B}$, si $\nu(B) = 0$ entonces $B=\varnothing$, y así $\mu(B) = \mu(\varnothing) = 0$. Se sigue $\mu \ll \nu$.
%----------------------------------------------------
\item % (ii)
Consideremos $A \en\mathcal{B}$ como cualquier conjunto con $\mu(A) > 0$ (por ejemplo $A = [0,1]$). Entonces, si $E \en\mathcal{B}$ es cualquier conjunto con $\nu(E) < \infty$, tenemos que $A\cap E$ es finito, por lo que $\mu(A \cap E) = 0$ (sabemos que los conjuntos contables tienen medida de Lebesgue nula); sin embargo, $\mu(A)>0$. Esto prueba que $\mu$ no es regular interior con respecto a $\nu$.
%----------------------------------------------------
\item % (iii)
Procedamos por contradicción suponiendo que existe tal $f$. Por un lado, es evidente que $f$ no es nula, pues si lo fuera, implicaría que $\mu$ es nula. Entonces existe $x\en X$ tal que $f(x) > 0$, así:
\[
\mu(\set{x})
= \int_{\set{x}} f \dd{\nu} = f(x)\nu(\set{x}) = f(x)>0
\]
No obstante, sabemos que $\mu(\set{x}) = 0$, por lo que esto es una contradicción. El error estuvo en suponer que existe tal $f$.
\end{enumerate}
\end{sol}